%\section{Вступ}
\section*{\centering ВСТУП}\addcontentsline{toc}{section}{ВСТУП}

%Застосування методів машинного навчання для генерації контенту - квал.
%Застосування багатошарових перцептронів для генерації текстового контенту - курс.

% \textbf{Обґрунтування вибору теми дослідження.} 

% \textit{Актуальнiсть дослiдження зумовлена стрiмким розвитком та зростанням популярностi методiв машинного навчання для генерацiї текстового контенту у останнi роки. Особливо значний iнтерес до цiєї галузi виник пiсля представлення у 2022 роцi чат-боту ChatGPT вiд OpenAI \cite{chatgpt}, який надав користувачам iнтерфейс на основi великої мовної моделi GPT для генерацiї текстiв природною мовою. Поява таких систем вiдкриває новi можливостi для автоматизацiї створення рiзноманiтного текстового контенту, вiд звичайних повiдомлень до технiчної документацiї та творчих текстiв. Водночас, стрiмкий розвиток цiєї сфери ставить новi виклики щодо оцiнки якостi, контролю та практичного застосування згенерованих текстiв. Тому актуальним є дослiдження сучасного стану розвитку методiв машинного навчання для генерацiї текстового контенту, аналiз нових пiдходiв, метрик оцiнювання, наборiв даних та сфер застосування, що з'явились у науковiй лiтературi протягом останнiх рокiв (2022-2024). Такий огляд дозволить узагальнити та систематизувати актуальнi здобутки у данiй галузi, окреслити перспективнi напрямки подальших дослiджень та розробок.}

\textbf{Обґрунтування вибору теми дослідження.} 
Стрімкий розвиток цифрових технологій та їх інтеграція в усі сфери життя суспільства зумовлюють необхідність формування у підростаючого покоління компетентностей з інформаційної безпеки. Сучасні старшокласники активно використовують цифрові пристрої та інтернет-сервіси, однак, як свідчать дослідження, їхні знання про кібербезпеку залишаються на недостатньому рівні \cite{Mishra202436}. Проблема навчання основ інформаційної безпеки учнів загальноосвітніх закладів набуває особливої актуальності в умовах зростання кіберзагроз та поширення цифрового шахрайства \cite{Spirin_Kovalchuk_2011}.

Реформування системи освіти в Україні, зокрема впровадження концепції Нової української школи \cite{nush_mon_compressed}, передбачає оновлення змісту навчання та застосування сучасних педагогічних технологій. Державний стандарт повної загальної середньої освіти \cite{KMU_898_2020} визначає інформаційно-комунікаційну компетентність як одну з ключових, що включає безпечне та відповідальне використання цифрових технологій. Модельні навчальні програми з інформатики \cite{inform_7_9_rivkind} передбачають вивчення питань інформаційної безпеки, однак обсяг навчального часу, відведеного на цю тематику, є обмеженим.

Водночас, наукові дослідження останніх років засвідчують необхідність розробки спеціалізованих курсів та модулів з кібербезпеки для учнів середньої школи \cite{Jerman2024, Ismail2024}. Ефективне навчання кібербезпеки потребує інноваційних підходів, що враховують вікові та індивідуальні особливості учнів \cite{Jerman20259093}. Зарубіжний досвід свідчить про успішне впровадження інтерактивних методів та практико-орієнтованих завдань у навчанні інформаційної безпеки \cite{Antunes2021, Wang2020}.

Диференційований підхід до навчання є одним із провідних напрямів сучасної педагогіки, що забезпечує можливість урахування індивідуальних особливостей учнів \cite{Tomlinson2001}. Українські науковці активно досліджують проблеми диференціації навчання в контексті реформування освіти \cite{ШАРОВ2021, Шевчук_2020}. Диференційоване навчання розглядається як комплексна система, що передбачає визначення складності та темпу навчання відповідно до індивідуальних особливостей учнів, застосування відповідних форм, методів і засобів навчання \cite{Пісоцька2022}. Реалізація диференційованого навчання в компетентнісній освіті сприяє підвищенню якості освітнього процесу \cite{Білик2017}.

Проведений аналіз наукової літератури та освітньої практики виявив низку суттєвих суперечностей:
\begin{itemize}
    \item між зростаючими вимогами суспільства до рівня сформованості компетентностей з інформаційної безпеки в учнів та реальним рівнем підготовки випускників ліцеїв;
    \item між потенціалом диференційованого навчання у формуванні компетентностей з інформаційної безпеки та недосконалістю існуючих методичних інструментів його реалізації;
    \item між необхідністю впровадження сучасних підходів до навчання інформаційної безпеки учнів ліцеїв та обмеженістю методичного забезпечення вибіркових модулів з даної тематики.
\end{itemize}

% Виявлені суперечності зумовлюють актуальність дослідження та визначають його проблему: які методичні засади диференційованого навчання вибіркового модулю ``Інформаційна безпека'' сприятимуть підвищенню якості освітнього процесу в ліцеях?
% Проведений аналіз наукової літератури та освітньої практики виявив низку суттєвих суперечностей [1-9]:
% \begin{itemize}
%     \item між та реальним рівнем підготовки випускників ліцеїв;
%     \item між потенціалом диференційованого навчання та недосконалістю існуючих методичних інструментів його реалізації;
%     \item між необхідністю впровадження диференційованого навчання та обмеженістю традиційних педагогічних технологій.
% \end{itemize}

\textbf{Об’єкт дослідження}~-- освітній процес у ліцеях.

\textbf{Предмет дослідження}~-- диференційоване навчання вибіркового модулю “Інформаційна безпека” учнів ліцеїв.

\textbf{Мета дослідження}~-- полягає в теоретичному обґрунтуванні та розробленні часткової методики диференційованого навчання вибіркового модулю “Інформаційна безпека” учнів ліцею.

Відповідно до мети визначено такі основні \textbf{завдання дослідження}: % для квал. роб.
\begin{enumerate} [label=\arabic*)]% дописать
    \item дослідити сучасні підходи до навчання інформаційної безпеки учнів ліцеїв;
    
    \item узагальнити та систематизувати теоретичні засади диференційованого навчання;

    \item розробити завдання для диференційованого навчання модулю ``Інформаційна безпека'' учнів ліцею;
    
    \item розробити методичні рекомендації щодо впровадження диференційованого навчання учнів ліцеїв під час вивчення вибіркового модуля ``Інформаційна безпека''. 
    
\end{enumerate}

\textbf{Методи дослідження}:
\begin{itemize}[label=--]
    \item аналіз науково-методичної літератури з проблем диференційованого  навчання учнів ліцеїв;
    \item систематизація педагогічного досвіду диференційованого навчання;
    \item теоретичне обґрунтування методичних підходів до  диференційованого навчання вибіркового модулю ``Інформаційна безпека''.
\end{itemize}


\textbf{Наукова новизна результатів дослідження} полягає в тому, що 
\begin{enumerate} [label=--]% дописать

    \item 
вперше обґрунтовано нові часткові методики диференційованого навчання  вибіркового модулю ``Інформаційна безпека'' учнів ліцеїв;
    \item 
подальшого розвитку набула система завдань з інформаційної безпеки для учнів ліцеїв;
    \item 
удосконалено підходи до диференційованого навчання учнів ліцеїв.

%    \item 
%проаналiзовано ефективність застосування запропонованих диференційованих завдань до вибіркового модулю ``Інформаційна безпека''.
\end{enumerate}



\textbf{Практичне значення результатів дослідження} полягає в тому, що 1) розроблено завдання для проведення занять з вибіркового модулю “Інформаційна безпека”, 2) створено методичні рекомендації до вибіркового модуля ``Інформаційна безпека'' для учнів ліцеїв. Запропоновані навчально-методичні матеріали можуть бути використані в освітньому процесі закладів загальної середньої та профільної освіти, а також здобувачами вищої педагогічної освіти під час вивчення методичних навчальних дисциплін, методистами та вчителями в системі підвищення кваліфікації та перепідготовки педагогічних працівників.


\textbf{Структура та обсяг роботи.} Робота складається зі
вступу, двох розділів, висновків до них, загальних висновків, списку використаних джерел (\total{citenum} найменувань), 4 додатка. Робота містить \totaltables\  таблиць та \totalfigures\ рисунки.
Загальний обсяг роботи~-- \pageref{last_page} 
сторінки.



% ABSTRACT
% ABSTRACT
% 2 • Report an abstract addressing each item in the PRISMA 2020 for Abstracts checklist.

% INTRODUCTION
% METHODS
% RESULTS
% DISCUSSION
% OTHER INFORMATION